% Context from chapters/sec-autodiff.tex; for mathematical reference, not compilation.
t dimension is $p=\sum_{k=1}^s n_k$.

\begin{figure}
\centering
\includegraphics[width=.7\linewidth]{auto-diff/comp-graph}
\caption{\label{fig-compgraph}
A computational graph.
}
\end{figure}

Represent the computation as a sequence of elementary maps:
\eq{
	\foralls k=s+1,\ldots,t, \quad x_k = f_k( x_1,\ldots,x_{k-1} )
}
Each $f_k$ depends only on earlier variables; see Figure~\ref{fig-compgraph}. The directed acyclic graph (DAG) connects each argument used by $f_k$ to its output $x_k$. The indexing is a topological order: every edge points from a smaller to a larger index.
 
Evaluating $f(x)$ is a forward traversal of the graph.
  
Automatic differentiation differentiates the supplied computational graph. It does not by itself choose an efficient implementation of the original formula; the derivative computation benefits from the efficiency and reuse already present in the forward program.



                                                                        
\subsection{Forward Mode of Automatic Differentiation}

Forward mode propagates derivatives of each intermediate variable with respect to a chosen input block, here $x_1$. For scalar inputs, repeating the traversal for all $p$ coordinates gives the full derivative. The formulas below also propagate a block Jacobian when $n_1>1$.

\newcommand{\bk}[1]{\left[#1\right]}
\newcommand{\pdb}[2]{\bk{\pd{#1}{#2}}}

\begin{figure}
\centering
\includegraphics[width=.4\linewidth]{auto-diff/forward-mode}
\quad
\includegraphics[width=.4\linewidth]{auto-diff/backward-mode}
\caption{\label{fig-forward-backward}
Derivative dependencies in forward mode (left) and reverse mode (right).
}
\end{figure}

Initialize the input derivatives by
\eq{
	\pd{x_1}{x_1} = \Id_{n_1 \times n_1}, \quad 
	\pd{x_2}{x_1} = 0_{n_2 \times n_1}, \ldots, \quad
	\pd{x_s}{x_1} = 0_{n_s \times n_1}, 
}
These are one and zero for scalar inputs. Propagate the derivatives using
\eq{
	\foralls k=s+1,\ldots,t, \quad
	\pd{x_k}{x_1} 
	= \sum_{\ell \in \text{parent}(k)} \bk{\pd{x_k}{x_\ell}} \times \pd{x_\ell}{x_1}
	= \sum_{\ell \in \text{parent}(k)} \pd{f_k}{x_\ell}(x_1,\ldots,x_{k-1}) \times \pd{x_\ell}{x_1}.
}
The notation ``$\text{parent}(k)$'' denotes the nodes $\ell<k$ of the graph that are connected to $k$, see Figure~\ref{fig-forward-backward}, left. 
 
The stored quantities are the derivatives $\pd{x_\ell}{x_1}$. The symbol $\times$ denotes matrix multiplication in the vector-valued case.
 
The brackets $[\ldots]$ indicate a local derivative: $\pd{x_k}{x_\ell}$ is the derivative of the elementary map $f_k$ with respect to that argument. It can be evaluated when needed, assuming the local derivative is available in closed form.

Suppose the local derivative operations $\pd{f_k}{x_\ell}$ have comparable costs, as when the variable dimensions $n_k$ are uniformly bounded, and each node has a bounded number of parents. Computing the full Jacobian coordinate by coordinate requires $p$ traversals, each with cost comparable to a forward evaluation. This becomes expensive for large $p$.



   
\paragraph{Simple example.}


We consider the function
\eql{\label{eq-simple-func-autodiff}
 	f(x,y) = y\log(x)+\sqrt{y\log(x)}
} 
on the open domain $x>0$, $y\log x>0$, where all elementary derivatives exist. Its computational graph is displayed in Figure~\ref{fig-dag-example-simple}. The iterations of the forward mode to compute the derivative with respect to $x$ read
\begin{align*}
		\pd{x}{x} &= 1,  \quad \pd{y}{x} = 0 \\
		\pd{a}{x} &= \bk{\pd{a}{x}} \pd{x}{x} = \frac{1}{x} \pd{x}{x} &
			\{x \mapsto a = \log(x)\}\\
		\pd{b}{x} &= \bk{\pd{b}{a}} \pd{a}{x}  + \bk{\pd{b}{y}} \pd{y}{x}  = y \pd{a}{x} + 0 &
			\{(y,a) \mapsto b=ya\}\\
		\pd{c}{x} &= \bk{\pd{c}{b}} \pd{b}{x}  = \frac{1}{2\sqrt{b}} \pd{b}{x}&
			\{b \mapsto c=\sqrt{b}\} \\
		\pd{f}{x} &= \bk{\pd{f}{b}} \pd{b}{x} + \bk{\pd{f}{c}} \pd{c}{x}  = 1\pd{b}{x} + 1\pd{c}{x}&
			\{(b,c) \mapsto f=b+c\}
\end{align*}
A second forward pass gives the derivative with respect to $y$:
\begin{align*}
		\pd{x}{y} &= 0,  \quad \pd{y}{y} = 1 \\
		\pd{a}{y} &= \bk{\pd{a}{x}} \pd{x}{y} = 0 &
			\{x \mapsto a = \log(x)\}\\
		\pd{b}{y} &= \bk{\pd{b}{a}} \pd{a}{y}  + \bk{\pd{b}{y}} \pd{y}{y}  = 0 + a \pd{y}{y}  &
			\{(y,a) \mapsto b=ya\}\\
		\pd{c}{y} &= \bk{\pd{c}{b}} \pd{b}{y}  = \frac{1}{2\sqrt{b}} \pd{b}{y}&
			\{b \mapsto c=\sqrt{b}\} \\
		\pd{f}{y} &= \bk{\pd{f}{b}} \pd{b}{y} + \bk{\pd{f}{c}} \pd{c}{y}  = 1\pd{b}{y} + 1\pd{c}{y}&
			\{(b,c) \mapsto f=b+c\}
\end{align*}



\begin{figure}
\centering
\includegraphics[width=.6\linewidth]{auto-diff/example-simple}
\caption{\label{fig-dag-example-simple}
Example of a simple computational graph.
}
\end{figure}

     
   
                           

                        
                                    
                               
  
                                                                                                                                                                                                   
              
                   
                                                                 
                              
                                                                         
                             
                                                                                               
                          
            



              
          
                                                               
                                        
                                              
 
            

   
                                 

                            
                                 
                                                                                      
 
                                                                                                                       
              
                   
                                                       
                           
                                                                       
                               
                                                                                                               
                                   
                                                                                                            
                                        
                                                      
                             
                                                                                                             
                                           
                          